% !TeX root = index.tex
% define document class
\documentclass{scrreprt}

% set variables
\newcommand{\varAuthor}{Max Muster}
\newcommand{\varCandidate}{\textbf{\varAuthor} \\ max.muster@example.com} % Kandidat/in
\newcommand{\varResponsibleSpecialist}{\textbf{Franziska Müller} \\ franziska.mueller@example.com} % verantwortliche Fachkraft
\newcommand{\varVocationalTrainer}{\textbf{Thea Meier} \\ thea.meier@example.com} % Berufsbildner/in
\newcommand{\varPrimaryExpert}{\textbf{Franz Krebs} \\ franz.krebs@example.com} % Hauptexperte
\newcommand{\varSecondaryExpert}{\textbf{Adrian Kuster} \\ adrian.kuster@example.com} % Nebenexperte
\newcommand{\varCompany}{Beispiel AG} % Firmenname
\newcommand{\varCompanyDepartment}{Entwicklung} % Abteilungsname
\newcommand{\varTitle}{Meine IPA}
\newcommand{\varVersion}{0.1} % Versionsnummer: Pro IPA-Tag um 0.1 erhöhen
\newcommand{\varExaminationBoard}{Prüfungskomission 19} % Prüfungsorganisation
\newcommand{\varExaminationBoardDepartment}{Informatik Applikationsentwicklung} % Fachrichtung

% apply ipa package
\usepackage{../lib/ipa}

\lohead[\varTitle]{\varTitle}
\lofoot[\today]{\today}
% see B6.5a
\cfoot[\varAuthor\ / \varCompany\\Version \varVersion]{\varAuthor\ / \varCompany\\Version 1}
\rofoot[Seite \pagemark{} von \pageref{LastPage}]

% make glossaries
\makenoidxglossaries

% define glossary entries
% see B6.6
\newglossaryentry{Organigramm}{
  name={Organigramm},
  description={Ein Organigramm stellt eine Organisation und deren Aufbauorganisation grafisch dar.}
}

% create document
\begin{document}

  % set page numbering
  \pagenumbering{roman}

  % set page style
  \pagestyle{scrheadings}

  % include title page
  \thispagestyle{empty}
  
  \begin{center}
  \makebox[\textwidth]{\includegraphics[width=1.1\paperwidth]{images/title.jpg}}
\end{center}
\vspace*{1cm}
{\fontsize{28}{28}\selectfont \varTitle}\\[0.25cm]
{\fontsize{16}{16}\selectfont \varAuthor}\\[0.25cm]
{\fontsize{16}{16}\selectfont \varCompany}

  \newpage
  \TileWallPaper{\paperwidth}{\paperheight}{images/background.pdf}

  % see A1.3
  % see B6.3a
  % generate the table of contents
  \tableofcontents

  % finish page
  \clearpage

  % use the arabic numbering system
  \pagenumbering{arabic}

  % reset page counter
  \setcounter{page}{1}

  % see B6.1a
  % create a phantom toc entry for "Umfeld und Ablauf"
  \clearpage\phantomsection\addcontentsline{toc}{part}{Umfeld und Ablauf}

  % [2]: Seite 11
\chapter{Aufgabenstellung}

Dieses Kapitel beinhaltet die komplette Aufgabenstellung im originalen Wortlaut.

\section{Ausgangslage}

\lipsum[1-2]

\section{Detaillierte Aufgabenstellung}

\lipsum[3-9]

\section{Mittel und Methoden}
\lipsum[10]
  % [2]: Seite 11
\chapter{Deklaration}

Folgender Abschnitt beschreibt die Vorkenntnisse des Kandidaten und dessen Vorbereitung.

\section{Vorkenntnisse}

\lipsum[11]

\section{Vorarbeiten}

\lipsum[12]

\section{Neue Lerninhalte}

\lipsum[13]

\section{Arbeiten in den letzten 6 Monaten}

\lipsum[14]

  % [1]: Kriterium B6
\chapter{Projektaufbauorganisation}

  % see B6.2b
\chapter{Zeitplan}
  % see B6.2c
\chapter{Arbeitsjournal}

  % see B6.1a
  % create a phantom toc entry for "Projekt"
  \clearpage\phantomsection\addcontentsline{toc}{part}{Projekt}

  % see A1.1a
\chapter{Informieren}
  % see A1.4
\chapter{Planen}
  \chapter{Entscheiden}
  \chapter{Realisieren}
  \chapter{Kontrollieren}
  \chapter{Auswerten}

  % see B6.6a
  % create a phantom toc entry for the index/glossary table
  \clearpage\phantomsection\addcontentsline{toc}{part}{Glossar}

  % generate glossary
  \printnoidxglossary[title={Glossar}]

  % create a phantom toc entry for the figures table
  \clearpage\phantomsection\addcontentsline{toc}{part}{Abbildungsverzeichnis}

  % generate figures table
  \listoffigures

  % defines the beginning of the appendix
  \appendix

  % create a phantom toc entry for "Projekt"
  \clearpage\phantomsection\addcontentsline{toc}{part}{Anhang}

  \chapter{Quellcode}

\end{document}
